This chapter highlights the implementiation of how options and modifiers are handled by the parser. 

The \verb|OptionModifierFormat| class is used to handle specific keywords. It includes two main functions, which are \verb|handle_option| and \verb|handle_modifier|. Each of these functions matches for its respective name. 

\section{Option @NO\_KETOP}

When encountered, the \verb|no_ketop| variable gets enabled in the \verb|current| \\(\verb|FormattedResponse|).

If that variable has already been set to true (i.e.\ if \verb|@NO_KETOP| was used twice in the file) a warning will be logged.

The \verb|@NO_KETOP| option expects no argument. 

\section{Modifier @VERSION}

TODO

\section{Modifier @OS [F]}

The \verb|@OS| modifier represents a conditional statement and requires an argument. 

The argument represents the operating system to check; it is compared against the \verb|os| variable provided in \verb|FormatterArgs|. This comparison is case-insensitive ignores any leading whitespace. 

If the provided \verb|os| variable is \verb|None|, a warning will be logged and the body gets discarded. 

If the comparison is satisfied, the body will be further recursively parsed; otherwise, the body will be discarded.

\section{Modifier @PREAMBLE [F]} 

The \verb|@PREAMBLE| modifier continues to recursively parse the body text, however the parsed text will be appended to the respective index of \verb|preamble_dict|, as opposed to \verb|script|. This index is given by the required argument.

This is done by returning a \verb|ChildArgs| instance with the \verb|preamble_index| variable set to that given by the argument. 

A syntax exception will be thrown if it is defined inside \verb|@BRACKETSTART|, \verb|@BRACKETEND| or another instance of \verb|PREAMBLE|.

\section{Modifiers @BRACKETSTART and @BRACKETEND [F]}

The \verb|@BRACKETSTART| or \verb|@BRACKETEND| modifier continue to recursively parse the body text, however the parsed text will be appended to \verb|bracketstart| or \verb|bracketend| respectively, as opposed to \verb|script|. 


This is done by returning a \verb|ChildArgs| instance with \verb|in_bracketstart| or \\\verb|in_bracketend| respectively set to \verb|True|.

A syntax exception will be thrown if it is defined inside \verb|@BRACKETSTART|, \verb|@BRACKETEND| or \verb|@PREAMBLE|.

The \verb|@BRACKETSTART| and \verb|@BRACKETEND| modifiers expect no argument.
